% 第 5 题三视图：按固定源内容流矢量层精确值复刻（1 pt → 0.03515 cm，与全库同制）。
% 原点取正/侧视图底边左端（源页 (145.4980, 675.4114) pt），y 向上；三视图共用 31.1815 pt/单位（足尺）。
% 源实测：正视图＝直角三角形（底 2、右边高 1，直角在右下）；侧视图＝等腰三角形（底 2 高 1）＋
% 顶点到底边中点的虚线（节距 on/off 2.989）；俯视图＝左边竖直（2）、右尖三角，无尺寸标注。
% 尺寸线：断口式（数字落在两段线之间）；尺寸界线各 6.236 长；箭头实心（本图用 stealth 三角同制）。
\begin{tikzpicture}[x=0.03515cm,y=0.03515cm,>=stealth,every node/.style={font=\small,inner sep=0pt},
  gl/.style={line width=0.398bp,line cap=round,line join=round}]
% ===== 正（主）视图：直角在右下 =====
\draw[gl] (0,0)--(62.3630,0)--(62.3630,31.1820)--cycle;
% 竖直尺寸 1（右边外侧）：尺寸界线 + 断口尺寸线 + 上下箭头
\draw[gl] (65.4817,0)--(71.7177,0);
\draw[gl] (65.4817,31.1817)--(71.7177,31.1817);
\draw[gl,->] (68.5995,9.3546)--(68.5995,0);
\draw[gl,->] (68.5995,21.8271)--(68.5995,31.1820);
\node at (68.5995,15.5906) {1};
% 水平尺寸 2（底边下方）：尺寸界线 + 断口尺寸线 + 左右箭头
\draw[gl] (0,-9.3543)--(0,-3.1183);
\draw[gl] (62.3630,-9.3543)--(62.3630,-3.1183);
\draw[gl,->] (24.9454,-6.2362)--(0,-6.2362);
\draw[gl,->] (37.4179,-6.2362)--(62.3630,-6.2362);
\node at (31.1820,-6.2362) {2};
% ===== 侧（左）视图：等腰三角形 + 顶点到底边中点虚线 =====
\draw[gl] (93.5450,0)--(155.9080,0)--(124.7270,31.1820)--cycle;
\draw[gl,dash pattern=on 2.989pt off 2.989pt] (124.7266,0)--(124.7266,31.1820);
% 底边两段 1（共用中间界线，箭头背对背）
\draw[gl] (93.5450,-9.3543)--(93.5450,-3.1183);
\draw[gl] (124.7266,-9.3543)--(124.7266,-3.1183);
\draw[gl] (155.9080,-9.3543)--(155.9080,-3.1183);
\draw[gl,->] (102.8996,-6.2362)--(93.5450,-6.2362);
\draw[gl,->] (115.3720,-6.2362)--(124.7266,-6.2362);
\draw[gl,->] (134.0812,-6.2362)--(124.7266,-6.2362);
\draw[gl,->] (146.5537,-6.2362)--(155.9080,-6.2362);
\node at (109.1360,-6.2362) {1};
\node at (140.3170,-6.2362) {1};
% ===== 俯视图：左边竖直、右尖三角（无尺寸标注）=====
\draw[gl] (0,-31.1819)--(0,-93.5449)--(62.3630,-62.3629)--cycle;
% ===== 三幅题注（照源字形中心定位；括号按全库先例用全角）=====
\node at (31.1820,-19.4786) {正（主）视图};
\node at (124.7270,-19.4786) {侧（左）视图};
\node at (31.1810,-100.6786) {俯视图};
\end{tikzpicture}
